\title{FlashAttention-3: Fast and Accurate Attention with Asynchrony and Low-precision}

\begin{abstract}
Within the widely adopted Transformer architecture, attention serves as a fundamental component but also represents a significant computational bottleneck, particularly in scaling large language models and supporting extended context lengths. \textsc{FlashAttention} previously introduced a strategy for accelerating attention computations on GPUs by significantly reducing memory access overheads. Nonetheless, this approach does not fully leverage advancements in the latest hardware; for instance, \textsc{FlashAttention-2} utilizes only 35\% of the H100 GPU’s computational potential. To address this, we propose three primary innovations tailored for Hopper GPUs: (1) employing warp specialization to utilize the asynchronicity of Tensor Cores and TMA, thus allowing concurrent data transfer and computation, (2) interleaving matmul and softmax computations at the block level, and (3) adopting block-level quantization and incoherent operations to make full use of the FP8 low-precision functionality built into modern hardware. Through these techniques, our new approach, \textsc{FlashAttention-3}, delivers 1.5-2.0$\times$ acceleration on H100 GPUs—achieving up to 740 TFLOPs/s using FP16 (representing 75\% utilization) and nearly 1.2 PFLOPs/s with FP8. Furthermore, we show that our FP8 \textsc{FlashAttention-3} implementation offers a 2.6$\times$ reduction in numerical error compared to standard FP8 attention baselines.
\end{abstract}

\section{Introduction}
\label{sec:intro}

Within the Transformer framework~\citep{vaswani2017attention}, the computation of self-attention presents a substantial performance constraint, as evaluating attention scores between queries and keys leads to quadratic computational growth relative to sequence length. Addressing attention efficiency in long-context scenarios promises to enable a range of new functionalities, including effective handling and inference across multiple extensive documents~\citep{guo2021longt5,shaham2022scrolls,peng2023yarn} and large-scale code repositories~\citep{roziere2023code, li2023starcoder}, as well as expansion to novel data types such as high-definition images~\citep{chen2022scaling}, audio~\citep{gulati2020conformer}, and video~\citep{ho2022video}. It also paves the way for new practical applications, for instance, systems capable of utilizing lengthy user interaction histories~\citep{sun2019bert4rec} or managing agents requiring extended operational planning~\citep{yao2022react}. These prospects have spurred a surge of research into accelerating attention mechanisms for long-context processing, through various avenues: deploying approximations~\citep{katharopoulos2020transformers,choromanski2020rethinking, tay2020efficient}, enhancing implementation efficiency~\citep{rabe2021self,dao2022flashattention,kwon2023efficient}, or proposing fundamentally new architectures~\citep{peng2023rwkv,sun2023retentive,gu2023mamba}.

This paper extends the advancements made by \citet{dao2022flashattention}, who developed exact-attention algorithms that carefully account for the GPU’s execution paradigms and architectural nuances within their broader algorithmic framework. Dao et al. \citep{dao2022flashattention} proposed \textsc{FlashAttention}, introducing a unique tiling methodology for attention computation that circumvents the need for intermediate data transfer to and from slow global memory by consolidating all attention-related operations into a unified GPU kernel. Subsequently, \citet{dao2023flashattention2} refined this approach with \textsc{FlashAttention-2}, which expands parallelism across the sequence length and restructures the forward pass’s inner loop to operate on blocks of the key and value matrices, thereby optimizing both resource usage and workload allocation on the GPU. Despite these improvements, our analysis reveals that \textsc{FlashAttention-2} demonstrates relatively low efficiency on contemporary GPU hardware when compared to specialized matrix multiplication (GEMM) kernels—for instance, realizing only 35\% utilization as opposed to the 80-90\% seen on Hopper H100 GPUs. This suboptimal performance appears to stem in part from differences at the implementation level, such as the failure to leverage Hopper-optimized instructions and instead utilizing those suited for Ampere architectures when accessing the Tensor Cores. Recent advancements, such as those presented in ThunkerKitten~\citep{spector2024thunder} and cuDNN 9~\citep{cudnn9}, demonstrate that employing Hopper-specific instructions in conjunction with tile-oriented abstractions can both accelerate attention mechanisms and streamline the codebase.

At a more foundational level, the algorithm implemented in \textsc{FlashAttention-2} is based on a straightforward synchronous model and does not deliberately integrate asynchronous operations or leverage low-precision computation in its core methodology. Asynchrony stems from advancements in hardware that accelerate vital operations within machine learning workflows; for instance, matrix multiplications are carried out by specialized components like Tensor Cores, and memory transfers are handled by units such as the Tensor Memory Accelerator (TMA), functioning independently from traditional CUDA cores, which are responsible for logic, integer, and floating-point computations. The adoption of low precision formats—including FP8 in Hopper, FP4 in Blackwell, and earlier formats like FP16 (introduced with Pascal in 2017) and BF16 (seen in Ampere in 2020)—has demonstrably allowed for a two- or fourfold increase in throughput without additional energy or area costs. We detail how Hopper supports these hardware features in \cref{subsec:hardware}. 
Redesigning \textsc{FlashAttention-2} to fully exploit these advancements involves resolving challenges such as orchestrating asynchronous execution between matmul and softmax (despite their interdependencies) and managing quantization carefully under low-precision settings, particularly in handling outlier values found in LLM features~\citep{dettmers2208llm, sun2024massive}.

In response, we introduce \textsc{FlashAttention-3}, a new algorithmic framework that incorporates and integrates three novel strategies to maximize performance on the latest GPU hardware architectures:\footnote{Our evaluations focus on NVIDIA’s Hopper architecture, though the methodology applies to any modern GPU architecture that effectively supports asynchronous execution and low-precision arithmetic.}

\begin{enumerate}

\item \textbf{Producer-Consumer asynchrony:} We introduce a warp-specialized approach to software pipelining that leverages asynchrony between data transfer operations and computations performed on Tensor Cores. This is accomplished by assigning data-producing and data-consuming tasks to distinct warps, thereby broadening the pipeline’s ability to mask memory and instruction issue delays.
\item \textbf{Hiding softmax under asynchronous block-wise GEMMs:} The relatively slow non-GEMM tasks involved in softmax, including floating point multiply-adds and exponentials, are executed concurrently with asynchronous WGMMA GEMM instructions. To facilitate this, we reformulate the \textsc{FlashAttention-2} algorithm, breaking up specific stepwise dependencies between softmax and GEMMs. For instance, in the two-stage approach, as one block of the score matrix undergoes softmax processing, a separate asynchronous WGMMA invocation proceeds to compute the next block.
\item \textbf{Hardware-accelerated low-precision GEMM:} We update the forward pass method so that FP8 Tensor Cores can be utilized for GEMM operations, resulting in nearly a twofold improvement in observed TFLOPs/s. This adaptation necessitates addressing the layout conventions for WGMMA, specifically with respect to storage patterns of FP32 accumulators and FP8 operand matrices. Techniques such as block quantization and incoherent processing are employed to reduce the accuracy degradation typically introduced by switching to FP8.
\end{enumerate}

To empirically assess our approach, we conduct a series of benchmarks of \textsc{FlashAttention-3} on an H100 SXM5 GPU across diverse parameter configurations. Our findings demonstrate that (1) in the forward pass, FP16 provides a 1.5–2.0$\times$ performance improvement over \textsc{FlashAttention-2} (peaking at 740 TFLOPs/s), while in the backward pass, speedups of 1.5–1.75$\times$ are observed; (2) in FP8 precision, the throughput approaches 1.2 PFLOPs/s; and (3) at longer sequence lengths, FP16 surpasses existing alternatives, with FP8 showing comparable performance\footnote{Specifically, with head dimension 64, \textsc{FlashAttention-3} FP8 takes the lead, whereas for head dimensions 128 and 256 it matches performance without causal masking, but falls behind when causal masking is applied.} to NVIDIA's cuDNN state-of-the-art attention implementation. We further confirm that FP16 \textsc{FlashAttention-3} exhibits numerical errors on par with \textsc{FlashAttention-2}, and outperforms conventional attention implementations owing to retention of intermediate values (such as softmax rescaling) in FP32. Additionally, under the use of block quantization and incoherent processing, FP8 \textsc{FlashAttention-3} delivers 2.6$\times$ greater accuracy compared to standard attention methods utilizing per-tensor quantization when confronted with outlier features.

We have released \textsc{FlashAttention-3} under a permissive license\footnote{\textsc{FlashAttention-3} is publicly available at \texttt{https://github.com/Dao-AILab/flash-attention}} and intend to incorporate it into PyTorch and Hugging Face ecosystems, aiming to make it accessible to a broad spectrum of practitioners and researchers.

\section{Background: Multi-Head Attention and GPU Characteristics}
\label{sec:background}

\subsection{Multi-Head Attention}
\label{subsec:multi_head_attn}

Suppose $\mathbf{Q}, \mathbf{K}, \mathbf{V} \in \mathbb{R}^{N \times d}$ denote the query, key, and value matrices for a single attention head, with $N$ representing the length of the sequence and $d$ corresponding to the dimensionality of each head. The resulting attention output, $\mathbf{O}$, is calculated as follows:

\begin{equation*}
  \mathbf{S} = \alpha \mathbf{Q} \mathbf{K}^\top \in \mathbb{R}^{N \times N}, \quad \mathbf{P} = \mathrm{softmax}(\mathbf{S}) \in \mathbb{R}^{N \times N}, \quad \mathbf{O} = \mathbf{P}\mathbf{V} \in \mathbb{R}^{N \times d},
\end{equation*}

Here, the $\mathrm{softmax}$ function operates on each row individually, and the scaling parameter $\alpha$ is commonly chosen as $1/\sqrt{d}$. To mitigate numerical instability issues that arise when using the exponential function, it is standard practice to subtract $\mathrm{rowmax}(\mathbf{S})$ from $\mathbf{S}$. In the context of multi-head attention (MHA), separate projections for queries, keys, and values are utilized for every head, enabling parallelized computation across all heads and batches, which results in the complete output tensor.

Consider a scalar loss function $\phi$, and denote the gradient operator as $\mathbf{d}(-) = \partial \phi / \partial (-)$. Given the gradient with respect to the outputs, $\mathbf{dO} \in \mathbb{R}^{N \times d}$, we derive the gradients with respect to $\mathbf{Q}$, $\mathbf{K}$, and $\mathbf{V}$ by applying the chain rule, resulting in the following:

\begin{align*}
  \mathbf{dV} &= \mathbf{P}^\top \mathbf{dO} \in \mathbb{R}^{N \times d} \\
  \mathbf{dP} &= \mathbf{dO} \mathbf{V}^\top \in \mathbb{R}^{N \times N} \\
  \mathbf{dS} &= \mathrm{dsoftmax} (\mathbf{dP}) \in \mathbb{R}^{N \times N} \\
  \mathbf{dQ} &= \alpha \mathbf{dS} \mathbf{K} \in \mathbb{R}^{N \times d} \\
  \mathbf{dK} &= \alpha \mathbf{dS}^\top \mathbf{Q} \in \mathbb{R}^{N \times d},
\end{align*}

Specifically, for the softmax derivative, if $p = \mathrm{softmax}(s)$, then $\mathbf{d}s$ is computed as $(\mathrm{diag}(p) - p p^\top)\mathbf{d}p$ for a given vector $s$. The notation $\mathrm{dsoftmax}(\mathbf{dP})$ represents this transformation applied row-wise across the input. Moreover, these calculations are fully amenable to parallelization both over the multiple attention heads and across batches in the backward computation of MHA.

\subsection{GPU hardware characteristics and execution model}
\label{subsec:hardware}

We describe the aspects of the GPU's execution model relevant for \textsc{FlashAttention-3}, with a focus on the NVIDIA Hopper architecture as a concrete instantiation of this model.

\paragraph{Memory hierarchy:} The GPU architecture features a stratified memory system, where higher bandwidth is associated with smaller storage capacities (\cref{tab:gpu-hierarchy})\footnote{\citet{luo2024benchmarking} provides a shared memory bandwidth figure of 128 bytes per clock cycle per SM, which is then scaled by 132 SMs and a boost clock frequency of 1830 MHz in our calculations.}. At the top, global memory (GMEM)—also referred to as HBM—serves as the off-chip DRAM shared by all streaming multiprocessors (SMs). Information residing in GMEM is automatically stored in the on-chip L2 cache for faster access. Further down the hierarchy, each SM possesses a limited but fast on-chip shared memory (SMEM), which is explicitly managed by the programmer and divided into many banks. The lowest tier is occupied by the register file present within each individual SM.

\paragraph{Thread hierarchy:} In the GPU programming paradigm, execution is structured through various logical groupings of threads. Starting from the smallest unit to the largest, the hierarchy consists of individual threads, warps (which comprise 32 threads), warpgroups (each formed by 4 consecutive warps), threadblocks (also referred to as cooperative thread arrays or CTAs), threadblock clusters (specific to Hopper architecture), and finally, grids.

The interrelationship between these two hierarchies is significant. Threads that belong to a single CTA are simultaneously scheduled onto the same SM, while CTAs residing within a cluster are jointly mapped to the same GPC. All threads inside a CTA have direct access to SMEM; in contrast, each individual thread possesses up to 256 dedicated registers (RMEM) for exclusive use.

\paragraph{Asynchrony and warp-specialization:}

GPUs function as throughput-oriented processors, leveraging parallelism and asynchronous operations to mask both memory and execution delays. Hopper introduces the Tensor Memory Accelerator (TMA), a specialized piece of hardware, to facilitate asynchronous data transfers between GMEM and SMEM \cite[\S7.29]{cuda}. In addition, differing from earlier architectures like Ampere, Hopper's Tensor Core—accessible using the warpgroup-wide WGMMA instruction \cite[\S9.7.14]{ptx}—also operates asynchronously and is capable of fetching input data straight from shared memory.

The incorporation of hardware mechanisms for asynchrony enables the development of warp-specialized kernels, wherein the warps within a CTA are designated strictly as either producers, responsible solely for data transfer, or consumers, dedicated exclusively to computation. This separation generally enhances the compiler's capability to produce highly efficient instruction pipelines \citep{warp-specialization-2011}. Moreover, Hopper architecture facilitates the flexible allocation of register resources across warpgroups by means of the \verb|setmaxnreg| directive \cite[\S9.7.17.1]{ptx}. This allows warps that perform matrix-multiply-accumulate (MMA) operations to utilize more RMEM than those conducting tensor memory access (TMA), where just a single thread suffices.

\paragraph{Low-precision number formats:}
\label{sec:low-precision-gpu}
Recent GPU designs feature dedicated hardware units aimed at boosting the performance of low-precision arithmetic. For instance, on Hopper, the WGMMA operation can exploit FP8 Tensor Cores to achieve double the SM-level throughput relative to FP16 or BF16 computations.

Nonetheless, the proper usage of FP8 WGMMA requires an appreciation of the operand layout limitations. Consider the scenario of performing a GEMM operation multiplying $A \times B^{\top}$, where $A$ is of size $M \times K$ and $B$ is of size $N \times K$. The operand $A$ or $B$ is considered \emph{mn-major} if it is stored contiguously along the outer $M$ or $N$ dimension, whereas it is deemed \emph{k-major} if its contiguity lies along the inner $K$-dimension. For the FP16 variant of WGMMA, either mn-major or k-major arrangements are acceptable for operands residing in SMEM. In contrast, FP8 WGMMA exclusively supports the k-major format. Additionally, in cases such as attention mechanisms where consecutive GEMM operations are fused within the same kernel, incompatibilities between the layout of FP32 accumulators and FP8 operands hinder the seamless chaining of dependent FP8 WGMMA instructions.

In the context of attention, these layout restrictions entail certain modifications to the design of an FP8 algorithm, which we describe in \cref{sec:algofp8}.

\subsection{Standard Attention and Flash Attention} As described by \citet{dao2022flashattention}, \textbf{standard attention} refers to the approach where attention computation on GPUs is performed by explicitly storing the intermediate matrices $\mathbf{S}$ and $\mathbf{P}$ in HBM. In contrast, the core innovation of \textsc{FlashAttention} was to utilize a localized softmax reduction, thereby circumventing the costly intermediate memory accesses, and to combine the attention computation into a single GPU kernel. This localized softmax operation is implemented in lines \ref{code-ws:softmax_start}-\ref{code-ws:softmax_end} within the consumer mainloop of \cref{alg:flash3_wgmma_ws_only}, and is further complemented by the block-wise rescaling of $\mathbf{O}$. A straightforward derivation demonstrating that this method indeed yields $\mathbf{O}$ is provided in \cite[\S 2.3.1]{dao2023flashattention2}.

\section{FlashAttention-3: Algorithm}
\label{sec:algo}

This section outlines the \textsc{FlashAttention-3} algorithm. We concentrate on the forward pass for clarity, while details of the backward pass can be found in~\cref{sec:algo_ws_bwd}. We begin by showing how warp-specialization, along with a circular SMEM buffer, is incorporated into the foundational \textsc{FlashAttention-2} algorithm. Next, we detail the approach for leveraging the asynchronous operations of WGMMA to implement a two-stage pipeline that overlaps the GEMM and softmax computations. Lastly, we discuss the necessary adjustments for supporting FP8, covering aspects of data layout compatibility as well as ensuring accuracy through block quantization and handling incoherent processing.

\subsection{Producer-Consumer asynchrony through warp-specialization and pingpong scheduling}
\label{sec:algo_ws}

\paragraph{Warp-specialization}
Similar to \textsc{FlashAttention-2}, the forward computation in \textsc{FlashAttention-3} exhibits high degrees of parallelism along the axes of batch size, number of attention heads, and query sequence length. Consequently, it is sufficient to describe the algorithm at the CTA level, where the algorithm processes a tile $\mathbf{Q}_i$ of the query matrix to generate the corresponding output tile $\mathbf{O}_i$. For clarity, we will first outline the warp-specialization approach utilizing a circular SMEM buffer, omitting for now the details related to GEMM-softmax overlap. With $d$ denoting the head dimension and $N$ the total sequence length, a specified query block size $B_r$ is selected to partition $\mathbf{Q}$ into $T_r = \lceil \frac{N}{B_r} \rceil$ blocks, labeled as $\mathbf{Q}_1, .., \mathbf{Q}_{T_r}$.

\begin{algorithm}[H]
    \caption{\small\label{alg:flash3_wgmma_ws_only}\textsc{FlashAttention-3} forward pass \textbf{without} intra-consumer overlapping -- CTA perspective}
    \begin{algorithmic}[1]
\REQUIRE Input matrices $\mathbf{Q}_i \in \mathbb{R}^{B_r \times d}$, $\mathbf{K}, \mathbf{V} \in \mathbb{R}^{N \times d}$ residing in HBM; key block size $B_c$ such that $T_c = \lceil \frac{N}{B_c} \rceil$.
\STATE Set up the pipeline object for controlling barrier synchronization with a $s$-stage circular SMEM buffer.
\IF {executing as the producer warpgroup}
\STATE Release a pre-specified number of registers.
\STATE Initiate transfer of $\mathbf{Q}_i$ from HBM to shared memory.
\STATE Once done, signal the consumer that $\mathbf{Q}_i$ has been loaded.
\FOR{$0 \le j < T_c$}
    \STATE Block until the $(j\,\%\,s)$th buffer stage has been processed by the consumer.
    \STATE Fetch blocks $\mathbf{K}_j, \mathbf{V}_j$ from HBM to shared memory at buffer index $(j\,\%\,s)$.
    \STATE On completion, signal to consumers that $\mathbf{K}_j, \mathbf{V}_j$ are available.
\ENDFOR
\ELSE 
\STATE Reserve a specific number of registers determined by the consumer warp count.
\STATE Locally initialize $\mathbf{O}_i = (0) \in \mathbb{R}^{B_r \times d}$, $\ell_i = (0) \in \mathbb{R}^{B_r}$, and $m_i = (-\infty) \in \mathbb{R}^{B_r}$.
\STATE Wait until $\mathbf{Q}_i$ becomes available in shared memory.
\FOR{$0 \le j < T_c$}
\STATE Stall until $\mathbf{K}_j$ becomes available in shared memory.
\STATE Execute $\mathbf{S}_i^{(j)} = \mathbf{Q}_i \mathbf{K}_j^T$ (SS-GEMM operation). Synchronize and await completion. \label{alg:ws_only_gemm1}
\STATE Record $m_i^{\mathrm{old}} = m_i$; update $m_i = \mathrm{max}(m_i^{\mathrm{old}}, \mathrm{rowmax}(\mathbf{S}_i^{(j)}))$. \label{code-ws:softmax_start}
\STATE Calculate $\widetilde{\mathbf{P}}_i^{(j)} = \mathrm{exp}(\mathbf{S}_i^{(j)} - m_i)$, then update $\ell_i = \mathrm{exp}(m_i^{\mathrm{old}} - m_i) \ell_i + \mathrm{rowsum}(\widetilde{\mathbf{P}}_i^{(j)})$. \label{code-ws:softmax_end}
\STATE Wait until $\mathbf{V}_j$ is loaded in shared memory.
\STATE Update $\mathbf{O}_i = \mathrm{diag}(\mathrm{exp}(m_i^{\mathrm{old}} - m_i))^{-1} \mathbf{O}_i + \widetilde{\mathbf{P}}_i^{(j)} \mathbf{V}_j$ (RS-GEMM). Synchronize and await commit. \label{alg:ws_only_gemm2}
\STATE Mark the $(j\,\%\,s)$th buffer segment as ready for reuse by the producer.
\ENDFOR
\STATE Perform $\mathbf{O}_i = \mathrm{diag}(\ell_i)^{-1} \mathbf{O}_i$ and $L_i = m_i + \log(\ell_i)$ for output normalization.
\STATE Write the resulting $\mathbf{O}_i$ and $L_i$ to HBM, storing them as the $i$th segment of $\mathbf{O}$ and $L$.
\ENDIF
\end{algorithmic}
\end{algorithm}

In our adaptation of \cref{alg:flash3_wgmma_ws_only} for Hopper, we employ \verb|setmaxnreg| to manage (de)allocation tasks, utilize TMA to load $\mathbf{Q}_i$ and $\{ \mathbf{K}_j, \mathbf{V}_j \}_{0 \leq j < T_c}$, and leverage WGMMA for conducting GEMM operations within the main loop on the consumer side. The SS or RS prefixes specify whether the initial operand is drawn from shared memory or the register file, respectively. When analyzing the control flow of \cref{alg:flash3_wgmma_ws_only}, it should be observed that TMA load instructions benefit from asynchronous execution, thus the issuance of these loads is not blocked by previous incomplete loads. Additionally, during the first $s$ passes of the producer’s main loop, there are no wait instructions, as the buffer is still undergoing initial population.

\paragraph{Pingpong scheduling} Due to the asynchronous execution capabilities of WGMMA and TMA, combined with the benefits of warp-specialization, it becomes feasible to concurrently process the softmax operation for one warpgroup while another warpgroup performs the GEMM task. This possibility is particularly relevant when considering the substantial throughput disparity between matmul and non-matmul operations on state-of-the-art hardware accelerators. For instance, the H100 SXM5 GPU boasts a matmul (FP16) throughput of 989 TFLOPS, but only 3.9 TFLOPS are allocated to special functions like exponentials\footnote{The CUDA programming guide indicates that each streaming multiprocessor (SM) can handle 16 special function operations per clock cycle. Calculating 16 operations $\times$ 132 SMs $\times$ 1830 MHz gives a throughput of 3.9 TFLOPS for these operations.}, which are crucial for softmax. Considering the FP16 attention forward pass with a 128-dimensional head, the number of FLOPS devoted to matmul exceeds those for exponentials by a factor of 512, whereas the throughput for exponentials lags by 256 times. Consequently, executing the exponential can consume up to half the cycles required for matmul. This imbalance becomes even more pronounced in the FP8 scenario, where matmul throughput is further enhanced while the exponential throughput remains static.

Because the exponential function is handled by a dedicated hardware component (the multi-function unit), optimal performance would be achieved by initiating the exponential computation concurrently with the matmul operations on the Tensor Cores. To facilitate this overlap, we employ synchronization primitives (\texttt{bar.sync} instructions), ensuring that the GEMM operations (specifically, GEMM1, which computes $\mathbf{P} \mathbf{V}$ within a given iteration, and GEMM0, which processes $\mathbf{Q} \mathbf{K}^\top$ for the succeeding iteration) for warpgroup 1 execute prior to those for warpgroup 2. This arrangement causes the softmax operation for warpgroup 1 to proceed in parallel with the GEMM executions in warpgroup 2. Subsequently, these functions alternate: warpgroup 2 performs the softmax while warpgroup 1 handles GEMMs—a strategy commonly described as "pingpong" scheduling. A visual representation of this process appears in~\cref{fig:pingpong_scheduling}. 

Although the actual implementation of this pingpong approach may exhibit less distinct concurrency than the schematic implies, empirical observations indicate that it generally yields a performance benefit (such as increasing throughput from 570 TFLOPS up to 620--640 TFLOPS for FP16 forward propagation, with a head dimension of 128 and sequence length 8192).

\paragraph{Attention variants} When implementing multi-query attention~\citep{shazeer2019fast} and grouped query attention~\citep{ainslie2023gqa}, we adopt the tensor indexing methodology from \textsc{FlashAttention-2} to ensure $\mathbf{K}$ and $\mathbf{V}$ are not redundantly stored in HBM.

\subsection{Intra-warpgroup overlapping GEMMs and softmax}

Even within one warpgroup, we can overlap some instructions in the softmax with some instructions in the GEMMs. We describe one technique to do so.

Within the attention algorithm, computations performed in the main (inner) loop exhibit inherent sequential dependencies, which restrict opportunities for intra-iteration parallelism. For instance, the (local) softmax operation (see lines \ref{code-ws:softmax_start} to \ref{code-ws:softmax_end}) depends on the intermediate result $\mathbf{S}_i^{(j)}$ produced by the initial GEMM, and the subsequent GEMM similarly takes $\widetilde{\mathbf{P}}_i^{(j)}$—the output of softmax—as an input. This structure necessitates the use of wait statements on lines \ref{alg:ws_only_gemm1} and \ref{alg:ws_only_gemm2} in \cref{alg:flash3_wgmma_ws_only}, which results in serialized execution between GEMM and softmax steps. Nonetheless, these execution dependencies may be alleviated by introducing additional register-based buffers to enable pipelining between loop iterations. Building on this observation, we present a two-stage\footnote{It is important to note that while the number of pipeline stages used for overlapping is upper-bounded by the circular SMEM buffer’s number of stages $s$, it does not necessarily coincide with it.} GEMM-softmax pipelining strategy:

\begin{algorithm}[H]
    \caption{\small\label{alg:flash3_wgmma}\textsc{FlashAttention-3} consumer warpgroup forward pass}
    \begin{algorithmic}[1]
      \REQUIRE  Matrices $\mathbf{Q}_i \in \mathbb{R}^{B_r \times d}$ and $\mathbf{K}, \mathbf{V} \in \mathbb{R}^{N \times d}$ located in HBM, a key block size $B_c$, with total blocks $T_c = \lceil \frac{N}{B_c} \rceil$.
      \STATE Allocate a predefined number of registers according to the current number of consumer warps.
      \STATE Set up on-chip buffers, initializing $\mathbf{O}_i = (0) \in \mathbb{R}^{B_r \times d}$, and initialize vectors $\ell_i, m_i$ as $(0)$ and $(-\infty)$ in $\mathbb{R}^{B_r}$.
      \STATE Await transfer of $\mathbf{Q}_i$ and the initial $\mathbf{K}_0$ block into shared memory.
      \STATE Perform WGMMA-based multiplication to obtain $\mathbf{S}_{\mathrm{cur}} = \mathbf{Q}_i \mathbf{K}_0^T$. Ensure the operation is committed, then await its completion.
      \STATE Free the first stage (index $0$) in the key matrix $\mathbf{K}$ buffer.
      \STATE Update $m_{i}$, compute $\tilde{\mathbf{P}}_{\mathrm{cur}}$ and adjust $\ell_{i}$ from $\mathbf{S}_{\mathrm{cur}}$; apply necessary rescaling to $\mathbf{O}_i$.
      \FOR{$1 \le j < T_c - 1$}
        \STATE Await the readiness of $\mathbf{K}_j$ in shared memory.
        \label{code:mainloop_start}
        \STATE Initiate WGMMA computation for $\mathbf{S}_{\mathrm{next}} = \mathbf{Q}_i \mathbf{K}_{j}^T$. Commit but continue asynchronously. \label{code:first_wgmma}
        \STATE Await $\mathbf{V}_{j-1}$ data block to reside in shared memory.
        \STATE Launch WGMMA to compute and accumulate: $\mathbf{O}_{i} = \mathbf{O}_{i} + \tilde{\mathbf{P}}_{\mathrm{cur}} \mathbf{V}_{j-1}$. Commit without immediate blocking. \label{code:second_wgmma}
        \STATE Await completion of $\mathbf{Q}_i \mathbf{K}_{j}^T$ WGMMA.
        \STATE Derive the new $m_{i}$, calculate $\tilde{\mathbf{P}}_{\mathrm{next}}$, and update $\ell_{i}$ utilizing the contents of $\mathbf{S}_{\mathrm{next}}$. \label{code:softmax}
        \STATE Await the completion of the operation $\tilde{\mathbf{P}}_{\mathrm{cur}} \mathbf{V}_{j-1}$, then perform any required scaling adjustments to $\mathbf{O}_i$.
        \STATE Release the $(j\,\%\,s)$th stage from the $\mathbf{K}$ buffer and the $(j-1\,\%\,s)$th from the $\mathbf{V}$ buffer.
        \STATE Assign the value of $\mathbf{S}_{\mathrm{next}}$ to $\mathbf{S}_{\mathrm{cur}}$.
        \label{code:mainloop_end}
      \ENDFOR \STATE Await full load of $\mathbf{V}_{T_c - 1}$ into shared memory.
      \STATE Compute the final output increment:
      $\mathbf{O}_{i} = \mathbf{O}_{i} + \tilde{\mathbf{P}}_{\mathrm{last}} \mathbf{V}_{T_c - 1}$, utilizing WGMMA, then commit and ensure completion.

\cref{alg:flash3_wgmma} serves as a substitute for the consumer phase outlined in \cref{alg:flash3_wgmma_ws_only}, thus establishing the full \textsc{FlashAttention-3} workflow for FP16 computations. Conceptually, we refer to WGMMA synonymously with asynchronous GEMM operations. During the primary loop (spanning lines \ref{code:mainloop_start} through \ref{code:mainloop_end}), execution of the second WGMMA in iteration $j$ (see line \ref{code:second_wgmma}) is performed concurrently with the softmax computation associated with the subsequent iteration $j+1$ (line \ref{code:softmax}).

Although the pipelined architecture presented above promises improvements in theoretical throughput, there are important real-world challenges to address:
\paragraph{Compiler reordering} The pseudocode assumes a specific sequence of execution, yet in practice, the compiler (NVCC) frequently reorganizes instruction order to achieve optimal performance. Such reordering can interfere with the designed alternation between WGMMA and non-WGMMA instructions, potentially undermining both correctness and the expected acceleration. Examination of the produced SASS code, however, indicates that the anticipated instruction overlap is preserved (see Section~\ref{sec:2-stage-sass}).
\paragraph{Register pressure} Avoiding register spilling is crucial for maintaining peak performance. Nevertheless, the 2-stage pipeline paradigm inherently introduces the need for additional registers to track intermediate computation and synchronization between stages. Notably, an extra $\mathbf{S}_{\mathrm{next}}$ must reside in registers, with an additional consumption proportional to $B_r \times B_c \times \text{sizeof}(\text{float})$ per threadblock. Such augmented register allocation can constrain the use of larger block sizes—another strategy often leveraged for efficiency—which are similarly intensive in register usage. In practice, an informed compromise should be reached through careful profiling.
\paragraph{3-stage pipelining} Building upon the 2-stage technique, we put forward a 3-stage variant intended to further interleave computation, overlapping the second WGMMA phase with the softmax. This modification potentially delivers better utilization of Tensor Cores, yet it intensifies register usage by necessitating storage for an additional pipelined stage. Consequently, optimizing between expanding tile size and deepening the pipeline becomes increasingly challenging. The full specification and performance analysis of the 3-stage approach are available in ~\cref{sec:3-stage}.

\subsection{Low-precision with FP8}
\label{sec:algofp8}

\textbf{Efficiency: layout transformations.} Executing the forward pass of \textsc{FlashAttention-3} using FP8 precision introduces unique complications regarding tensor layout consistency that are not present when using FP16.

Typically, input tensors $\mathbf{Q}$, $\mathbf{K}$, and $\mathbf{V}$ are laid out contiguously along the head dimension. However, due to the k-major ordering required by FP8 WGMMA in the second GEMM, $\mathbf{V}$—specifically, the sub-blocks of $\mathbf{V}$ that are loaded into SMEM—must be organized contiguously along the sequence length dimension. Since the TMA load operation preserves the existing contiguous dimension, two primary approaches are possible: (1) performing a full transpose of $\mathbf{V}$ in GMEM prior to computation, or (2) transposing sub-blocks of $\mathbf{V}$ in-kernel after they have been brought into SMEM. 

For the first strategy, one can either (1a) fuse the transpose with the epilogue of the prior computation step, such as rotary embedding, or (1b) invoke a dedicated transpose kernel as a pre-processing step\footnote{Well-tuned transpose kernels can operate close to the peak bandwidth of the hardware~\citep{colfax_cutlass_transpose_2024}.} to swap the sequence length and head dimension strides. Nonetheless, option (1a) proves challenging to incorporate into general-purpose libraries, whereas option (1b) is inefficient in memory-constrained scenarios like inference due to additional memory traffic.

In the case of FP8 \textsc{FlashAttention-3}, we select option (2) as our approach. Within the kernel, the transpose operation leverages LDSM (\verb|ldmatrix|) and STSM (\verb|stmatrix|) instructions, which coordinate a warp’s threads to move data between shared memory (SMEM) and register memory (RMEM) in 128-byte chunks.\footnote{According to the PTX specification, LDSM/STSM are intended to transfer $8 \times 8$ matrices of 16-bit elements \cite[\S 9.7.13.4.15-16]{ptx}. For FP8, we circumvent this by packing pairs of 8-bit values to utilize these instructions. However, the transposed LDSM/STSM variants do not support unpacking 8-bit pairs, requiring extra register operations between LDSM and STSM to complete a tile-level transpose, though we exclude these specific details here.} These instructions are efficient in terms of register usage, which makes it feasible to carry them out in the producer warpgroup, and they also support transposed memory layouts during data transfer. Additionally, starting from the second iteration, we can schedule the transpose of the subsequent $\mathbf{V}$ tile to overlap with the two WGMMAs computations associated with the previous $\mathbf{V}$ tile and the current $\mathbf{K}$ tile, effectively hiding the transpose latency.

Second, it should be noted that, in contrast to FP16, the memory arrangement of the FP32 accumulator utilized in an FP8 WGMMA deviates from that assumed for operand A when stored in registers. Illustrations of partial layouts for both can be found in \cref{fig:rmem_accum} and \cref{fig:rmem_operand}, which display how data entries are allocated in registers for each thread according to their ordering. By leveraging byte permutation instructions, the accumulator from the initial WGMMA can be restructured to align with the format required by the subsequent WGMMA, while also ensuring compatibility with the $\mathbf{V}$ tile layout generated by the in-kernel transpose. Specifically, referring to \cref{fig:rmem_accum}, the data ordering is systematically changed to
$$\{ \verb|d0 d1 d4 d5 d2 d3 d6 d7| \},$$
with this pattern repeated for every 8-byte segment. Logically, with respect to the shape of the $\mathbf{P}$ tile, this action reorganizes its columns (e.g., columns $0189$ become the initial four columns). To facilitate WGMMA in generating the correct result tile, the in-kernel transpose may correspondingly produce a row permutation for the $\mathbf{V}$ tile that matches this column reordering.\footnote{The added flexibility offered by performing the transpose within the kernel removes the necessity of using shuffle instructions to transfer register contents between threads, as previously discussed in~\citep{colfax_fp8_flashattention_2024}.}

\textbf{Accuracy: block quantization and incoherent processing.} The FP8 (e4m3) representation allocates just 3 bits for the mantissa and 4 bits for the exponent, leading to greater numerical imprecision in comparison with FP16/BF16 formats. This increased error is further exacerbated in large-scale models, where outlier values~\citep{dettmers2208llm,
  sun2024massive} can differ substantially in magnitude from the majority of other elements, thereby presenting challenges for effective quantization. Common practice involves using per-tensor scaling~\citep{micikevicius2022fp8}, where a single scaling factor is assigned per tensor (such as one each for $\mathbf{Q}$, $\mathbf{K}$, and $\mathbf{V}$). To mitigate the accumulation of numerical inaccuracies when applying FP8 to attention mechanisms, we incorporate two particular strategies:

\begin{enumerate}

\item \textbf{Block quantization}: A single scaling factor is retained for each block; for $\mathbf{Q}$, $\mathbf{K}$, and $\mathbf{V}$, the tensors are partitioned into sub-blocks of dimension $B_r \times d$ or $B_c \times d$, and quantization is performed on these individual blocks. This quantization procedure may be integrated with an operation immediately preceding the attention mechanism, such as rotary embedding, without incurring additional computational latency, as rotary embedding is typically limited by memory bandwidth. Given that the \textsc{FlashAttention-3} algorithm processes data in blocks by design, the scaling for block quantization can be incorporated into each $\mathbf{S}$ block without any computational penalty.

\item \textbf{Incoherent processing}: To address the presence of outliers, $\mathbf{Q}$ and $\mathbf{K}$ are each pre-multiplied by a random orthogonal matrix $\mathbf{M}$ before FP8 quantization. Owing to the orthogonality of $\mathbf{M}$, where $\mathbf{M}\mathbf{M}^\top = I$, it holds that $(\mathbf{Q}\mathbf{M})(\mathbf{K}\mathbf{M})^\top = \mathbf{Q}\mathbf{K}^\top$; that is, pre-multiplying both $\mathbf{Q}$ and $\mathbf{K}$ by $\mathbf{M}$ preserves the attention computation. This strategy distributes outlier values across multiple components, since each entry in $\mathbf{Q}\mathbf{M}$ or $\mathbf{K}\mathbf{M}$ results from a random linear combination of the respective inputs, thereby diminishing the impact of quantization errors. Following the approach of \citet{chee2024quip} and \citet{tseng2024quip}, the matrix $\mathbf{M}$ is constructed as the product of random diagonal matrices of $\pm 1$ and a Hadamard matrix, allowing matrix multiplication to be performed with time complexity $O(d \log d)$ rather than $O(d^2)$, while still permitting fusion with the rotary embedding step at no extra computational overhead.

We validate that these two techniques reduces numerical error by up to 2.6$\times$ in \cref{sec:numerical_error}.

\section{Empirical Validation}
\label{sec:exp}

We use the primitives from CUTLASS~\citep{Thakkar_CUTLASS_2023} such as WGMMA and TMA abstractions to implement \textsc{FlashAttention-3} and evaluate its efficiency and accuracy.

\begin{itemize}

\item \textbf{Benchmarking attention.} We evaluate the performance of \textsc{FlashAttention-3} in terms of runtime for varying sequence lengths and perform a comparative analysis against the baseline PyTorch implementation, \textsc{FlashAttention-2}, and its Triton variant utilizing H100-optimized instructions, as well as the cuDNN-based vendor version of \textsc{FlashAttention-2} tailored for H100 GPUs. The results indicate that \textsc{FlashAttention-3} exhibits up to a 2.0$\times$ speed advantage over \textsc{FlashAttention-2} and is 1.5$\times$ faster relative to the Triton-optimized counterpart. Peak throughput for \textsc{FlashAttention-3} is observed at 740 TFLOPs/s, representing 75\% of the H100 GPU's theoretical peak performance.
\item \textbf{Ablation study.} Through systematic evaluation, we demonstrate that the enhancements provided by warp-specialization and the integration of GEMM-softmax pipelining significantly drive the increased performance in \textsc{FlashAttention-3}.
\item \textbf{Accuracy of FP8 attention.} Our analysis verifies that the adoption of block quantization and incoherent processing strategies leads to a 2.6$\times$ reduction in numerical error for FP8 computations within \textsc{FlashAttention-3}.
\end{itemize}

\subsection{Benchmarking Attention}
\label{subsec:benchmark_attn}

We evaluate the execution times of various attention mechanisms on an H100 80GB SXM5 GPU across multiple configurations (with or without causal masking, head dimension set to 64 or 128) using FP16 input data. The performance outcomes are presented in~\cref{fig:benchmark_attn_fwd} and~\cref{fig:benchmark_attn_bwd}. The data indicates that \textsc{FlashAttention-3} achieves approximately 1.5 to 2.0$\times$ speedup over \textsc{FlashAttention-2} during the forward computation and is about 1.5 to 1.75$\times$ faster in the backward computation. In comparison with conventional attention methods, \textsc{FlashAttention-3} demonstrates acceleration factors as high as 3 to 16$\times$. Notably, for sequence lengths of 1,000 tokens or more, \textsc{FlashAttention-3} can outperform even the proprietary, highly-optimized vendor library (cuDNN) specifically designed for H100 GPUs.

\paragraph{Benchmark settings:} Sequence lengths are adjusted from 512 up to 16k, increasing incrementally, while the batch size is chosen to ensure a constant total of 16k tokens. The hidden layer dimension is fixed at 2048. Head dimension values are tested at 64, 128, or 256, corresponding respectively to 32, 16, and 8 heads. For the purpose of measuring the forward pass FLOPs, we employ the following calculation:

\begin{equation*}
  4 \cdot \text{seqlen}^2 \cdot \text{head dimension} \cdot \text{number of heads}.
\end{equation*}

Applying causal masking requires dividing the total by 2, as only about half of the possible computations are actually performed. To estimate the FLOPs for the backward pass, we scale the forward pass FLOPs by a factor of 2.5. This is because the forward computation involves 2 matrix multiplications, whereas the backward pass (with recomputation) involves 5.

We additionally assess the forward pass runtime using FP8 precision under analogous experimental conditions. Results corresponding to headdim 256 can be found in~\cref{fig:benchmark_attn_fp8}, with the complete set of outcomes detailed in \cref{sec:benchmark_attn_fp8_full}.

\subsection{Ablation Study: 2-Stage Pipelining Experiments}
\label{sec:2-stage-pipelining-experiments}

We conduct an ablation study of both the 2-stage WGMMA-softmax pipelining and the warp-specialization techniques applied to the non-causal FP16 \textsc{FlashAttention-3}, maintaining fixed configuration parameters $\{\text{batch}, \text{seqlen}, \text{nheads}, \text{hdim}\} = \{ 4, 8448, 16, 128\}$. The findings presented in~\cref{table:ablation_pipelining} demonstrate that these algorithmic strategies—specifically, asynchrony through warp-specialization and the concurrency between GEMM and softmax operations—produce a marked performance improvement, boosting throughput from 570 to 661 TFLOPs.

\subsection{Numerical Error Validation}
\label{sec:numerical_error}

Given growing attention to the issue of numerical error in \textsc{FlashAttention}~\citep{golden2024flash}, we evaluate and contrast the precision of \textsc{FlashAttention-2}, \textsc{FlashAttention-3}, and a conventional attention implementation relative to an FP64 reference. To model the presence of extreme features and activations observed in LLMs~\citep{dettmers2208llm, sun2024massive}, we construct the entries of $\mathbf{Q}, \mathbf{K}, \mathbf{V}$ using the following distribution:

\begin{equation*}
  \mathcal{N}(0, 1) + \mathcal{N}(0, 100) \cdot \mathrm{Bernoulli}(0.001).
\end{equation*}

Specifically, every entry follows a normal distribution with a mean of zero and a standard deviation of 1, but for 0.1\% of the elements, an additional independent Gaussian noise with a standard deviation of 10 is included. The root mean squared error (RMSE) is subsequently evaluated and presented in~\cref{table:numerical_error}. When using FP16, both \textsc{FlashAttention-2} and \textsc{FlashAttention-3} yield RMSE values that are 1.7$\times$ smaller than those from the baseline method, as these algorithms retain intermediate results (namely, the softmax outputs) in FP32 precision. The FP8 baseline implementation utilizes per-tensor scaling, carries out matrix multiplication accumulation in FP32, and stores intermediate softmax values in FP16. Owing to the use of block quantization and incoherent processing, \textsc{FlashAttention-3} operating in FP8 demonstrates 2.6$\times$ higher accuracy than the FP8 baseline.

\section{Dicussion, Limitations, Conclusion}
\label{sec:discussion}

Through the development of \textsc{FlashAttention-3}, we illustrate that leveraging emerging programming strategies alongside recent hardware innovations, including asynchrony and low-precision computation, can significantly improve both the efficiency and precision of attention mechanisms. Our approach achieves a 1.5-2.0$\times$ speedup over \textsc{FlashAttention-2} and results in a 2.6$\times$ decrease in FP8 numerical error relative to standard per-tensor quantization methods. Nonetheless, certain aspects remain to be improved in future work: these include advancing optimization for LLM inference, incorporating a persistent kernel architecture into the FP8 kernel,\footnote{In our evaluations, the FP16 \textsc{FlashAttention-3} implementation utilizes a persistent kernel and a load balancing strategy, whereas the FP8 version currently does not. This discrepancy partly accounts for the inferior performance of FP8 \textsc{FlashAttention-3} at lower sequence lengths and under causal masking, especially when compared to the FP8 cuDNN kernels.} and further investigation into the impact of low-precision attention during extensive model training. Although our experiments have focused on Hopper GPUs, we anticipate that the presented methodologies will extend to a broad range of hardware accelerators. We further hope that the availability of a more rapid and precise attention primitive can enable new developments in long-context task domains.

\end{document}